\section{Finite history integrals}
\label{sec:histories}
\label{sec:history-estimates}

All estimates in this section concern a finite prescribed term of
Section~\ref{sec:cell}.  Its inputs have independent envelopes.  They do
not concern the probability of a good event for the interacting particle
process.  In particular, no estimate on an evolved factorial marginal is
used here.

\subsection{Sources, lifetimes, and a common energy weight}

Fix the discrete choices in a history: its labels, cell diagrams, tensor
factors, and retained slots at layer cuts.  Write $\Gamma$ for this
prescription and $\mathcal D_\Gamma$ for the domain of all its indicators.
The number of labels is $p$ and the number of final, or restarting, roots
is $h_0\geq1$.  Empty scalar terms have their usual separate interpretation.
Throughout this section $P=1+|\log\varepsilon|$.
Every label has one source time and one departure time.  A label at a
source is supplied by one independent factor; a label which survives to
the top departs at the observation time.  These times are fixed mesh
times.  The extra labels in a cell depart at its later endpoint in
forward time.  A contact time used below is an integration coordinate,
not a time at which the prescribed force system is changed.

We also include the positive terminal witness of
Definition~\ref{def:cell-terminal-witness}.  At its last cell the
distinguished roots are spatial anchors and normalization labels only.
There is no endpoint isolation requirement between those anchors and
the other final participants.  All labels of each final force group
depart together, and the whole terminal Hamiltonian and all of its
velocity coordinates are retained until integration.  Earlier cuts
obey the usual isolation rule.  Thus this variant does not cut a
persisting cluster into an anchor and its companions.  The number
$h_0$ still counts distinguished anchors; all $p-h_0$ other terminal
positions are integrated by the pins below.

Let the absolute value of the source for label $i$ be bounded by
\begin{equation}
 |u_i(x,v)|\leq \eta_i g_i(x)e^{-\beta_0|v|^2},\qquad
 \sum_{k\in\mathbb Z^3}\sup_{|x-k|\leq1}g_i(x)\leq M,
 \qquad 0\leq\eta_i\leq1 .
 \label{eq:hist-source-envelope}
\end{equation}
The constants $\beta_0>0$ and $M$ are fixed independently of the number
of layers.  The factors $\eta_i$ record small initial or residual
inputs.  Put $\eta_\Gamma=\prod_i\eta_i$.
One can take $g_i$ to be the Gaussian-weighted velocity supremum of
$u_i/\eta_i$ when $\eta_i>0$; if $\eta_i=0$, the term is zero. Thus
\eqref{eq:hist-source-envelope} follows from the norm used in
Propositions~\ref{pre:envelopes} and~\ref{cent:residual}.

Before the contact changes of variables, the absolute integral for
$\Gamma$ is the integral of the product of its absolute source factors over
$\mathcal D_\Gamma$ in the $p$ source phase coordinates, multiplied by
the inherited factorial coefficients and the activity factor
$\mu^{p-h_0}$. A test of the top phase points with absolute value at most
one may be included; those phase points are integrated as well.
An initial centered input is first replaced by its pointwise product
bound. We never divide this integral by the mass of $\mathcal D_\Gamma$.

For each label retain a position and a velocity coordinate at all times
in $[0,T]$.  Before its source and after its departure freeze both
coordinates.  On its lifetime evolve it with its prescribed force group.
At a fixed time this gives $p$ phase coordinates, including the frozen
ones.  Denote them by $\widehat Z(t)$, and write
\begin{equation}
 S=\sum_{i=1}^p|v_i^{\rm source}|^2,
 \qquad \mathcal E=S+2Bp.
 \label{eq:hist-source-energy}
\end{equation}
The quantity $S$ is a function of the whole history coordinates.  It is
not, in general, a quadratic function of the departure velocities.

\begin{lemma}[Augmented volume and the once-counted Gaussian]
\label{hist:gaussian}
For every fixed prescription $\Gamma$, the map from its $p$ source
coordinates to $\widehat Z(t)$ is one-to-one and volume preserving
outside null sets.  On $\mathcal D_\Gamma$,
\begin{equation}
 \sum_{i=1}^p|\widehat v_i(t)|^2\leq S+2Bp,
 \qquad
 \sup_{0\leq t\leq T}\sum_{\substack{i:\,i\text{ active at }t}}|v_i(t)|^2
       \leq S+2Bp.
 \label{eq:hist-mixed-energy}
\end{equation}
In particular, with $V(t)=(\widehat v_i(t))_{i=1}^p$,
\begin{equation}
 e^{-\beta_0 S}
 \leq e^{\beta_0Bp}
       e^{-\beta_0|V(t)|^2/2}e^{-\beta_0 S/2}.
 \label{eq:hist-split-gaussian}
\end{equation}
The same assertions hold at the last time, when the mixed coordinates
consist of the top coordinates and the departure coordinates of every
other label.  The exponent in \eqref{eq:hist-split-gaussian} is not
changed at a layer cut.
\end{lemma}

\begin{proof}
Between two consecutive mesh times the augmented map is a product of
Hamiltonian maps and identity maps.  Each factor is invertible and
volume preserving by Proposition~\ref{pre:flow}.  Injection merely changes
which of the already present coordinates is active; departure changes
which coordinate is frozen.  Both operations are identity maps on the
augmented space.  A fixed prescription therefore gives a finite
composition of volume preserving maps.  Coincident singular positions
at a switch form a null set; apply the argument first on compact sets
avoiding them and then exhaust the domain.

For the energy identity, keep the Hamiltonian of every force group.
At the earlier endpoint of a cell its distinct force groups are unions
of distinct whole initial polymers.  Their cross potentials are zero.
Thus passing from the preceding force partition to this partition
changes no interaction energy.  At the later endpoint the root group
$A$ splits into its surviving labels $R$ and $A\setminus R$.
The indicator $e_I^R(A)$ makes their cross potentials zero.  The other
force groups depart in their entirety.  Store the Hamiltonian of each
departing group at that time.  At a layer cut the singleton sources
are inserted on the restored polymer pattern, so their cross
potentials with each other and with all existing groups in the new
force system are zero.  A nonsingleton polymer is never split by this
operation.

At the top of a terminal witness all final groups are stored intact,
including their distinguished anchors.  No split into $R$ and
$A\setminus R$ is made there, and no endpoint isolation is used.
The mixed tuple still has one terminal velocity for every label.
This proves the same identity for the positive witness variant.

There is a relevant distinction at a late cell endpoint.  A detached
group may geometrically overlap a different-generation group or the
root group.  Those groups had separate Hamiltonians throughout the
cell.  Their geometric overlap introduces no term into either
Hamiltonian, and they are not merged when the detached group departs.
Consequently it causes no missing cross-energy term in this calculation.

Starting with the empty active system, conservation on open cells and
the preceding switch identities give, at every time $t$,
\begin{equation}
 \frac12\sum_{i:\,\text{source}_i\leq t}|v_i^{\rm source}|^2
   =\sum_{\substack{A:\,A\text{ current at }t}}\mathscr H_A(t)
       +\sum_{\substack{D:\,D\text{ departed by }t}}\mathscr H_D(\text{departure}_D).
 \label{eq:hist-energy-ledger}
\end{equation}
The current and departed groups partition the injected labels.  Stability
on each group therefore subtracts at most $B$ times the number of those
labels, not $B$ times that number at every cell.  Adding the kinetic
energy of future sources to both sides proves
\eqref{eq:hist-mixed-energy}.  Splitting the source exponential into
two equal factors and using the first inequality gives
\eqref{eq:hist-split-gaussian}.
\end{proof}

\begin{remark}[Positive potential energy is retained]
\label{hist:positive-potential}
The reverse inequality $S\leq |V(T)|^2+Cp$ is not asserted.  An
interacting group can depart with large positive potential energy and
small kinetic energy.  Bounds on intermediate speeds consequently
use $\mathcal E=S+2Bp$.  The last factor in
\eqref{eq:hist-split-gaussian} pays for functions of this source energy.
This distinction is necessary for singular repulsive cores and is
also useful for bounded potentials.
\end{remark}

Here is a convenient precise way of using that last factor.  Divide
the domain into the shells
\begin{equation}
 rp\leq S<(r+1)p,\qquad r=0,1,\ldots,\qquad
 E_r=(r+1+2B)p.
 \label{eq:hist-shells}
\end{equation}
On such a shell every active or mixed squared-velocity sum is at most
$E_r$.  After reserving any fixed positive part of the last exponential
in \eqref{eq:hist-split-gaussian}, the shell has an extra factor
$e^{-c rp}$.  Parent choices may now be summed using the same scalar
$E_r$ on every branch.  This is also a rigorous implementation of
expressions of the form
\begin{equation}
 C^p\sup_{s\geq0}\{W(s)e^{-cs}\}.
 \label{eq:hist-energy-payment}
\end{equation}
One must not first replace $S$ by a terminal kinetic energy, or pull its
branch-dependent value through a parent sum.

\subsection{Translation charts on a cell}

The following construction is used in two versions.  The fine version
preserves the original tensor-factor parent of a singleton addition;
it is used for exponential counting.  The coarse version is permitted
to record every geometric contact and is used when an extensive power
of $\varepsilon$ pays its larger combinatorial cost.  Neither version
alters a force in the prescribed history.

\begin{lemma}[First-contact translation charts]
\label{hist:pinning}
Consider one cell of length $\Delta$ with $q$ prescribed terminal root
labels and $m$ extra labels.  Products of simultaneous cell diagrams
are allowed.  Keep all root coordinates and all terminal velocities
fixed.  On the domain where every extra label is linked by the
prescribed cell contacts to a root, the extra terminal positions are
covered by finitely many charts with the following properties.

\begin{enumerate}
\item Each extra label costs one three-dimensional pin.  A pin is
either a relative endpoint displacement in $B(0,\varepsilon)$, with
Jacobian one, or a contact time $t$, a normal $\omega\in\mathbb S^2$,
and the Jacobian
\begin{equation}
 \varepsilon^2
 \bigl|(v_i(t)-v_j(t))\cdot\omega\bigr|\,dt\,d\omega.
 \label{eq:hist-contact-jacobian}
\end{equation}
The velocities in this formula are those of the separately computed
translation components immediately before their merger.
\item The number of discrete pin charts, including a cell's generation
choices, is at most $(C(p+1))^{Cm}$.  No partition of the roots is
introduced.  The full spatial Jacobian is the product of the pin
Jacobians.  No norm of a derivative of a many-particle flow occurs.
\item If $m=1$ globally in that cell and the extra label is separated
from its own roots at the late endpoint, its first backward contact
with its own tensor factor gives the exact chart
\begin{equation}
 y=x_a(t)+\varepsilon\omega+(b_I-t)w,
 \qquad
 dy\,dw=\varepsilon^2|(w-v_a(t))\cdot\omega|
          \dd t\,\dd\omega\,\dd w .
 \label{eq:hist-single-chart}
\end{equation}
Here $b_I$ is the late endpoint, $y,w$ are the extra terminal position
and velocity, and the root trajectory is computed without that extra
label.  Apart from the full or virtual choice in
Remark~\ref{rem:cell-single}, there is no multiplicity independent of
the choice of $a$.
\end{enumerate}
On the energy shell \eqref{eq:hist-shells}, multiplication by
$\mu^m$ and integration over the $m$ pin coordinates gives the coarse
bound
\begin{equation}
 (Cp)^{Cm}(1+E_r)^{m/2}(\Delta+\varepsilon)^m.
 \label{eq:hist-multi-cell}
\end{equation}
The convention $p+1$ can replace $p$ in this and subsequent bounds at
zero size.
\end{lemma}

\begin{proof}
At the late endpoint initially tie all root labels into one anchored
translation component.  This is a convention about fixed coordinates,
not an assertion that the roots interact.  Process cross-component
endpoint overlaps, choosing a pair by a fixed label order.  An overlap
fixes the relative translation by an offset of length at most
$\varepsilon$.  Ties between two already anchored pieces fix no new
coordinate and are ignored.  Repeat until every remaining floating
component is separated from all other components at that endpoint.

Every component carries the original force partition restricted to its
labels.  Translate the component as a whole.  Its restricted dynamics
commutes with that translation.  Starting at the late endpoint, evolve
all these components backward.  Until their first intercomponent
geometric contact, they exert no cross force and these restricted
motions equal the original prescribed motions.  At that contact fix
one relative component translation, merge the two translation
components, and continue with the restricted original force partition
on their union.  Physical interactions are included when the original
partition includes them; a virtual contact only merges translation
coordinates.  Interactions already internal to a component need not
be enumerated.  Because the prescribed contact graph links every
extra label to roots, the algorithm eventually anchors every extra
coordinate.

For two floating components write their paths as
$a_A+q_i^A(t)$ and $a_B+q_j^B(t)$, where all internal coordinates are
fixed.  At the selected contact,
\[
 a_B-a_A=\varepsilon\omega+q_i^A(t)-q_j^B(t).
\]
The derivative in time is the relative velocity; the other two
derivatives span the tangent plane of the sphere of radius
$\varepsilon$.  Their determinant is
\eqref{eq:hist-contact-jacobian}.  A contact with the anchored component
uses the same formula with its position already fixed.  Endpoint pins
are linear changes from a relative translation to an offset and have
determinant one.

At each stage keep the previously chosen internal and pin coordinates
unchanged.  The next change acts only on one still free relative
translation.  Thus the changes are triangular in merger order.  Large
derivatives with respect to earlier pin coordinates enter only
off-diagonal blocks.  This proves the product assertion, including
for singular potentials by compact exhaustion of the regular flow
domains.

There are at most $m$ mergers that remove an extra translation degree
of freedom.  Choosing their pairs, representatives, orders, and
endpoint/contact types costs at most $(Cp)^{Cm}$.  The extra-label
partitions and generations have the same type of bound by
Proposition~\ref{cell:expansion}.  Roots have already been anchored
together and require no partition sum.  On each original first-contact
domain the construction recovers the starting configuration.  For an
upper bound the area formula permits summing the images of all these
charts.  It does not require a uniform bound on repeated physical
contacts.  Tangencies lie in zero-Jacobian portions of the contact
maps, and deterministic tie rules suffice at multiple choices.

When there is one extra label, before its first contact with its own
factor that label is free and the old roots follow their isolated
root flow.  This proves \eqref{eq:hist-single-chart}.  An earlier
contact with a different tensor factor exerts no force and hence does
not invalidate this original-factor chart.  Its indicators remain
part of the term.  The force group and original factor are determined
by the chosen original parent and the full/virtual choice.

Finally, on the valid energy shell every velocity at a selected contact
has magnitude at most $E_r^{1/2}$.  A contact pin therefore integrates
to at most $C\varepsilon^2 E_r^{1/2}\Delta$, whereas an endpoint pin
costs $C\varepsilon^3$.  After $\mu^m=\varepsilon^{-2m}$,
\eqref{eq:hist-multi-cell} follows.  The energy bound is used on the
original chart domain before other restrictions are dropped; no
energy bound on an artificially reconstructed off-domain trajectory
is presumed.
\end{proof}

\begin{remark}[Coarse global pins and relative variables]
\label{hist:coarse-chart}
For a geometric estimate one may run the same algorithm with all
currently displayed root paths, across tensor factors, in the anchored
component.  A singleton then has its primary pin at its earliest
backward contact with any displayed path.  If this is in a different
factor, the primary contact is virtual and its later required
original-factor contact is still an indicator.  This is only a new
coordinate chart for the same force prescription.

At a singleton pin put $g=w-v_a(t)$.  All root trajectories at $t$
are determined before this new block is introduced.  The map
$w\mapsto g$ has determinant one.  Combining it with the normal-to-disk
map gives
\begin{equation}
 \mu\,dy\,dw=|g|\,dt\,dg\,d^2\zeta,
 \qquad \zeta\in g^\perp,\quad |\zeta|<1,
 \label{eq:hist-disk-chart}
\end{equation}
on the incoming hemisphere.  This identity is a contact-flux identity,
not an inverse differential cross-section.  A measurable orthonormal
frame identifies each disk with a fixed unit disk and preserves area.

These substitutions are triangular in backward cell order.  A whole
multiple-addition cell is one coordinate block; all its variables are
introduced before any earlier cell is computed.  Consequently an
ordinary parent state is determined by preceding blocks, never by
the fresh velocity or impact variable at that pin.  This remains true
with arbitrarily complicated old-root dynamics.

On $S\leq Rp$, after fixing source summability boxes, all velocity,
root-position, and remaining relative-coordinate ranges have
polynomial size in $p$.  Rescale endpoint offsets by $\varepsilon$.
The preceding determinant calculation then bounds the activity-normalized
chart density, in the product of contact blocks and remaining auxiliary
variables, by $(Cp)^{Cp}$ for a fixed cell prescription.  Summing
discrete cell prescriptions replaces this by $P^{C_{A,Q,L}p}$ when
$p\leq P^A$ and $J=\lceil P^Q\rceil$.  This coarse bound is used only
with an extensive geometric gain.  Reassigning a primary pin across
factors can lose the original factor choice; it is not used in the
exponential parent count below.  Section~\ref{sec:exposure} records the
additional exclusions that make the coarse charts causal for selected
terminal trees.
\end{remark}

\subsection{Spatial anchors and the finite exponential bound}

There is one independent three-dimensional velocity coordinate for
each formal label: the terminal velocity of a top or the departure
velocity of an extra.  Lemma~\ref{hist:pinning} changes position
coordinates only, before the optional relative substitution.  Thus
these $3p$ velocity coordinates, denoted by $V$, have a common Gaussian
majorant in all the original-factor charts.

For spatial integration use the cubes
$Q_k=k+[-1/2,1/2)^3$ and numbers
$a_{i,k}=\sup_{|x-k|\leq1}g_i(x)$.  Expand each spatial source envelope
over these cubes.  Their coefficients sum to at most $M$ per label.
For a fixed list of boxes, a top label and its unique source position
satisfy
\begin{equation}
 \sum_{\substack{i:\,i\text{ a top}}}|x_i^{\rm top}-x_i^{\rm source}|^2
 \leq T\int_0^T\sum_{\substack{i:\,i\text{ a top, active at }t}}|v_i(t)|^2dt
 \leq T^2E_r.
 \label{eq:hist-top-displacement}
\end{equation}
There is no positional jump at a change of force group.  After all
extra positions have been pinned, the remaining top positions
therefore range in a ball in $\mathbb R^{3h_0}$ of squared radius at
most $C_T(E_r+h_0)$ about their source boxes.  Its volume is bounded by
\begin{equation}
 \mathsf S(h_0,E_r)
   =\left[C_T\left(1+\frac{E_r}{h_0}\right)^{3/2}\right]^{h_0}.
 \label{eq:hist-spatial-factor}
\end{equation}
The other spatial box indicators may be discarded.  Summing all box
coefficients costs $M^p$.  In particular, for $h_0=o(p)$ and
$E_r=O(p)$ the logarithm of \eqref{eq:hist-spatial-factor} is $o(p)$.

Here and below cell occupancy means the total occupancy across all
simultaneous tensor factors.  Let $m_*$ count labels in cells with at
least two additions, and let $o_*$ count singleton endpoint pins.
The remaining $n_*=p-h_0-m_*-o_*$ additions are called ordinary pins.
The interlayer retained-slot choices are either fixed, or retained
explicitly as their binomial selection factors.  They are never
replaced by $2^{CH}$.

\begin{proposition}[Finite-term bound with its time factors]
\label{hist:basic}
Fix $A<\infty$ and the finite number $L$ of layers.  There is a
numerical exponent $d$ in the cell estimates such that, if
$J=\lceil P^Q\rceil$ with $Q>2dA$, all the following bounds hold
uniformly for $1\leq p\leq P^A$, for sufficiently small $\varepsilon$.
Put $n_j$ for the total additions in layer $j$, $h_j$ for its incoming
roots, and $q_j=h_j+n_j$.  After summing cell choices, cell diagrams,
and original parents, the absolute integral on the shell
\eqref{eq:hist-shells} is bounded by
\begin{align}
 &\eta_\Gamma C^p e^{-c rp}\mathsf S(h_0,E_r)
       \chi_\varepsilon^{m_*+o_*}(r+2)^{(m_*+o_*)/2}
       \prod_{j=1}^L\frac{b^{n_j}q_j^{n_j}}{n_j!}
 \nonumber\\[-2pt]
 &\hspace{22mm}\times
  \int_{\mathbb R^{3p}}
       \prod_{\substack{i:\,i\text{ an ordinary pin}}}
          \left(1+|w_i|^2+\frac{E_r}{q_{j(i)}}\right)^{1/2}
       d\gamma(V),
 \label{eq:hist-finite-bound}
\end{align}
where $d\gamma$ is a normalized product Gaussian of one fixed exponent,
the integral may retain the common indicator $|V|^2\leq E_r$,
and one may take
\begin{equation}
 \chi_\varepsilon
   = C P^{dA}\left(J^{-1/2}+\frac{\varepsilon J}{b}\right).
 \label{eq:hist-mesh-gain}
\end{equation}
Harmless increases of $d$ absorb the fixed extra powers occurring in
the finite diagrams.  Constants $C,c$ in
\eqref{eq:hist-finite-bound} do not depend on $L$ or on later rarity
cutoffs.  The threshold of sufficiently small $\varepsilon$ may
depend on the fixed $b=T/L$.

In particular, after summing layer occupancies and retained slots the
total finite integral of order $p$ is at most
\begin{equation}
 \eta_\Gamma C_L^p.
 \label{eq:hist-coarse-exponential}
\end{equation}
Its portion with $S>Rp$ is at most
\begin{equation}
 \eta_\Gamma C_L^p e^{-cRp},
 \label{eq:hist-source-tail}
\end{equation}
after changing $c>0$ and taking $R$ above a fixed constant.
These are estimates of finite independently enveloped terms, not
assertions about the actual process.
\end{proposition}

\begin{proof}
For a set of at most $q$ old labels with squared-velocity sum at most
$E$, direct Cauchy--Schwarz gives
\begin{equation}
 \sum_a|v_a-w|^2\leq2q|w|^2+2E,\qquad
 \sum_a|v_a-w|
       \leq Cq\left(1+|w|^2+\frac E q\right)^{1/2}.
 \label{eq:hist-parent-sums}
\end{equation}
The root state at an ordinary pin is determined before its parent and
new velocity are chosen.  On a fixed source-energy shell, retain the
adapted gate that this root state has squared-velocity sum at most
$E_r$.  Every original term in the shell satisfies that gate.  Sum
the last parent first using \eqref{eq:hist-parent-sums}, then the
preceding one, and continue.  The resulting factors are independent
of the parent states.  This proves the ordinary-pin product in
\eqref{eq:hist-finite-bound}.  Remaining future admissibility indicators
can now be dropped.  The majorant in $V$ is the fixed Gaussian of
Lemma~\ref{hist:gaussian}, so these integrations introduce no additional
per-layer exponent loss.

For globally ordinary cells the first contact times are chronologically
ordered.  Summing their disjoint cell boxes is an integral over a
subset of the corresponding ordered simplex.  Hence it supplies
$b^{n}/n!$, not $J^n$ times this quantity.  A maximal stretch of empty
cells simply composes the same prescribed root flows.  It has one
deterministic map and no independent empty-cell choice.

If a cell has $m\geq2$ additions, use
\eqref{eq:hist-multi-cell}, without a factorial improvement in its
internal contact times.  For a total of $m_*$ such participants there
are at most $m_*/2$ occupied cells.  Their location choices cost at
most $J^{m_*/2}$ in a layer.  The time factors supply
$(b/J)^{m_*}$, since $\varepsilon J\leq b$ eventually.  Distributing
the participants among those cells and choosing their pin/generation
data costs $p^{C m_*}$.  A missing ordered-time factorial for these
participants is restored with
$n_j!/(n_j-m_j)!\leq p^{m_j}$.  On the energy shell, their speed
factors cost at most $p^{m_*/2}(r+2)^{m_*/2}$ times a fixed exponential.
All these polynomial powers are absorbed in $d$ in
\eqref{eq:hist-mesh-gain}.

A singleton endpoint pin has factor $C\varepsilon$ after its activity
normalization, instead of a contact-time integral.  Its cell choice
costs at most $J$.  Restoring the time simplex and any parent factor
loses at most a further fixed power of $p$ per such pin.  This gives
the second term in \eqref{eq:hist-mesh-gain}.  The unordered-extra
factorials of Corollary~\ref{cor:cell-factorial} cancel permutations of
new names; at a singleton no independent factor assignment remains
after choosing its original parent.  At multiple cells all remaining
assignments are among the polynomial costs just counted.

The energy and spatial bounds were proved in
\eqref{eq:hist-split-gaussian}--\eqref{eq:hist-spatial-factor}.
Combining them with these changes of variables proves
\eqref{eq:hist-finite-bound}.  Notice that $\chi_\varepsilon\to0$
uniformly for $p\leq P^A$.  No rarity cutoff has entered its exponent.

For the coarse consequence replace every parent capacity by $p$ and
use the global ordered simplex.  For $n=p-h_0\leq p$,
\[
 p^n\frac{T^n}{n!}\leq e^p\max(1,T)^p.
\]
The source shell gives $\sum_i|w_i|^2\leq E_r$ on its original
domain.  Alternatively integrating the product Gaussian moments gives
the same bound $C^p(r+2)^{Cp}$ for the speed product and the spatial
factor, since $h_0\leq p$ and
$h_0\log(1+p/h_0)\leq Cp$.  Thus the shell sum is bounded by
\[
 C^p\sum_{r\geq0}e^{-c rp}(r+2)^{Cp}\leq C^p.
\]
All retained-slot binomials are at most
$2^{\sum_jq_j}\leq2^{(L+1)p}$, and all other fixed layer choices
cost a further $C_L^p$.  This proves
\eqref{eq:hist-coarse-exponential}.  Starting the shell sum at
$r\geq R-1$ proves \eqref{eq:hist-source-tail}.  Polynomial costs of
special cells have already been paid by their mesh factors before
this shell sum or a later constant-cutoff estimate is used.
\end{proof}

\subsection{Compact parameters and adapted parent counting}

For this subsection work on $S\leq Rp$ and put $R_1=R+2B$.
The original-factor charts have an explicit auxiliary product measure.
Normalize their Gaussian velocities and disk areas, choose every
ordinary parent uniformly from a fixed list of $p$ slots, and kill
inactive or otherwise illegal choices by indicators.  Root coordinates
and all auxiliary multiple-cell blocks are integrated against the
majorants already used above.  A whole such block is revealed before
proceeding to an earlier ordinary cell.  Fix the ordered contact times
when discussing a conditional mark.  The past of pin $i$ contains all
preceding blocks, its time, and its old-root state, but not its parent
choice or its new $w_i,\zeta_i$.

Let $G_i$ be the adapted gate that the old-root squared-velocity sum is
at most $R_1p$, together with any other adapted restrictions explicitly
retained.  Write $F_i=|v_{i,I_i}-w_i|$, where $I_i$ is the uniform
parent slot.  From \eqref{eq:hist-parent-sums},
\begin{equation}
 \mathbb E_{\rm ref}(G_iF_i^2\mid\text{past of }i)
       \leq C(1+R_1).
 \label{eq:hist-reference-square}
\end{equation}
For any prescribed set of $k$ pins whose marks have conditional
reference probability at most $\eta$, repeated conditional integration
gives a probability at most $\eta^k$.  Applying one Cauchy--Schwarz
inequality to the whole product of fluxes, and integrating the last
fresh block first in each factor, proves
\begin{equation}
 \mathbb E_{\rm ref}\!\left[
       \prod_iG_iF_i\,\mathbf1_{\{N_{\rm mark}\geq k\}}\right]
 \leq [C(1+R_1)]^{n_*/2}
             \binom{n_*}{k}^{1/2}\eta^{k/2}.
 \label{eq:hist-reference-discard}
\end{equation}
Physical parent velocities need not be independent.  What is used is
that they are measurable before their own new block is introduced.
Later admissibility restrictions may be dropped; this does not allow
one to drop a future dependence of a parent state.

\begin{proposition}[Finite parameter discards]
\label{hist:irregular}
Fix $L$, $A$, $\delta>0$, and $M_0<\infty$.  The source-energy cutoff
$R$ and fixed numbers
\[
 U,W<\infty,\qquad \gamma,\kappa>0
\]
can be chosen so that, for all sufficiently small $\varepsilon$,
the total finite history integral with at least $\delta p$ marked
additions is at most $\eta_\Gamma e^{-M_0p}$.  An addition is marked
if it is in a multiple-addition cell, is an endpoint pin, or is an
ordinary pin satisfying any of the following conditions:
\begin{enumerate}
\item its original parent speed exceeds $U$, its added absolute
velocity exceeds $W$, or its relative speed is below $\gamma$;
\item for a genuine non-force-free isolated continuation, either
output increment relative to the incoming parent has magnitude below
$\kappa$;
\item its isolated two-body duration in physical time exceeds
$\varepsilon^{.96}$;
\item its contact time is within $\varepsilon^{.95}$ of a cell
boundary, or it belongs to a pair of birth times separated by less
than $2\varepsilon^{.9}$.
\end{enumerate}
The contribution of $S>Rp$ has the same bound.  The constants in this
proposition may depend on $L$; they do not enter the $C$ in
\eqref{eq:hist-finite-bound}.
\end{proposition}

\begin{proof}
Choose $R$ using \eqref{eq:hist-source-tail}.  Given an old-root state
with squared-velocity sum at most $R_1p$, at most $R_1p/U^2$ of its
parent slots have speed above $U$.  The conditional reference
probability of choosing one is therefore at most $R_1/U^2$.  A large
added absolute velocity has Gaussian probability at most
$Ce^{-cW^2}$.  The ball $|w-v|<\gamma$ has reference probability at
most $C\gamma^3$, uniformly in $v$.

On the remaining compact set let $G_*=U+W$.  In the incoming disk
variables, relative velocities and impact parameters have reference
density bounded by a constant with respect to $dg\,d^2\zeta$.
On $\gamma\leq|g|\leq G_*$ this measure is equivalent to incoming
flux.  Proposition~\ref{pre:scattering} therefore gives a function
$\omega_\Phi(\kappa;\gamma,G_*)\to0$ as $\kappa\downarrow0$ such
that the second mark has conditional probability at most
$C\omega_\Phi$.  The bound is uniform over $|v|\leq U$: translate
$w$ to $g=w-v$, bound its Gaussian density above, and enlarge to the
fixed relative-velocity annulus.  Force-free passages are crossings
and are not marked for their zero parent increment.

Lemma~\ref{pre:duration} gives
\[
 \int_{|g|\leq G_*}\int_{|\zeta|<1}
       |g|\,\mathbf1_{\{\tau(g,\zeta)>D\}}\,d^2\zeta\,dg
       \leq C_{\Phi,G_*}/D.
\]
Dividing by $|g|\geq\gamma$ shows that the duration mark has
conditional probability at most
$C_{\Phi,G_*}\gamma^{-1}\varepsilon^{.04}$.
This is the duration of the isolated continuation, not an assumption
that an actual many-particle encounter is isolated.  An interruption
will be an encounter defect in Section~\ref{sec:exposure}.

Take $U,W$ large and then $\gamma,\kappa$ small to make the sum of
the first two reference-mark probabilities as small as needed.
Equation~\eqref{eq:hist-reference-discard},
$\binom{n_*}{k}\leq2^p$, and the finite exponential normalization
give the required bound for an extensive number of these marks.
For fixed $\delta$ any prescribed exponential is obtained by making
their one-step probability sufficiently small.  Multiple cells and
endpoint pins have instead the factor
$\chi_\varepsilon^{m_*+o_*}$ from
Proposition~\ref{hist:basic}, which proves their assertion.

For the time marks use unordered uniform times before sorting in
each layer.  A fixed time lies within $\varepsilon^{.95}$ of a cell
boundary with probability at most
$CJ\varepsilon^{.95}/b$.  If $k$ birth times have a neighbor at
distance below $2\varepsilon^{.9}$, their closeness graph on the
line has a matching or a disjoint consecutive-pair selection of
size at least a fixed multiple of $k$.  For a prescribed disjoint
pair selection the probability is at most
$(C_L\varepsilon^{.9})^{ck}$.  There are at most $p^{Ck}$ selections.
Times in distinct layers which are this close lie near their common
boundary and can be charged there.  Since $p\leq P^A$ and $J$ is
a power of $P$, these bounds tend to zero with a positive power of
$\varepsilon$ per fixed fraction of marked times.  Apply the same
global Cauchy--Schwarz estimate to the fluxes.  A finite union bound,
with thresholds reduced by the number of categories, completes the
proof.
\end{proof}

\subsection{A two-sided volume estimate for crowding}

Crowding here concerns the displayed paths of a finite term.  In
particular, they may belong to different prescribed force groups.
The frozen extension is useful: its positions are continuous at every
fixed lifetime and force-partition switch, and on $S\leq Rp$ their
speeds are bounded by $C_R\sqrt p$.

\begin{lemma}[A fixed-pivot forest]
\label{hist:pivot}
Fix a time $t_*$, a forest of $r$ edges on the $p$ labels, and
$\rho\geq\varepsilon$.  For one fixed force/lifetime prescription
whose required contacts connect every label to a top, the source integral
before multiplication by $\mu^{p-h_0}$, on $S\leq Rp$ with
$|x_i(t_*)-x_j(t_*)|\leq\rho$ on every forest edge is at most
\begin{equation}
 \eta_\Gamma C^p p^{Cp}\rho^{3r}
       \varepsilon^{2(p-r-h_0)_+}.
 \label{eq:hist-pivot-bound}
\end{equation}
More precisely, a stratum with $k$ final translation components has
the factor $\rho^{3r}\varepsilon^{2(p-r-k)}$, with $k\leq h_0$.
There is no factor equal to the number of cells in this geometric
bound.
\end{lemma}

\begin{proof}
Use the augmented phase coordinates at $t_*$ from
Lemma~\ref{hist:gaussian} and first fix their velocities.  Root each
component of the forest at its least label.  Replace its positions
by one common translation and the edge differences along that rooted
tree.  The linear map has determinant one.  The $r$ relative vectors
range in balls of radius $\rho$, at cost $(C\rho^3)^r$.

Each current translation component carries the original time-dependent
force partitions restricted to its labels, as well as the frozen
coordinates outside their lifetimes.  Begin by merging any remaining
cross-component overlaps at the pivot, using endpoint offsets in
$B(0,\varepsilon)$.  Next scan \emph{all} geometric contacts of the
extended paths, both forward and backward from $t_*$.  This includes
contacts between different physical force systems and contacts of
inactive frozen coordinates.  Choose the first remaining contact in
increasing distance $|t-t_*|$, with a deterministic rule for ties,
fix its relative translation, and merge the two translation components.
Recompute the merged component's restricted prescribed motion from
the pivot in both directions.  Continue until no cross-component
contact remains anywhere in $[0,T]$.

The reconstruction has a useful induction.  Inside the already
explored two-sided time window, distinct translation components have
had no geometric contact.  Hence they exert no physical force, and
their separately computed motions agree with the full prescribed
motion.  A newly chosen earliest contact lies on the boundary of an
extension of that window.  Merging its two components changes no
motion inside the old window.  This statement holds even when
successive mergers occur on opposite sides of the pivot.  Interactions
beyond the known window are recomputed with the original force
partitions after the merger.  A virtual merger never adds a force.

The contact Jacobian is again
\eqref{eq:hist-contact-jacobian}; for a frozen coordinate its position
derivative is zero.  Its integral over all possible times is bounded
by $C_T\varepsilon^2\sqrt{R_1p}$.  Piecewise smooth time intervals
are summed as one time integral.  There is no independent endpoint
choice at each mesh time: extended positions are continuous, so an
overlap which becomes physically relevant at a switch was either
already a pivot overlap or was preceded by a geometric sphere
contact.  This is why all extended geometric contacts, rather than
only currently interacting pairs, were included in the scan.

The triangular determinant and area-formula argument of
Lemma~\ref{hist:pinning} applies verbatim.  At a fixed-time kink use
the adjacent smooth time intervals; the endpoint itself has zero
three-dimensional contact-parameter measure unless it is covered by
an already counted overlap.  Tangent and repeated contacts need no
separate counting.  Compact exhaustion avoids singular physical
configurations.  There are at most $p-1$ mergers, with at most
$p^{Cp}$ choices of pairs, directions, orders, and contact/overlap
types.

After the first $r$ forest links, if the final number of components
is $k$, there have been exactly $p-r-k$ further pins.  At termination
there are no remaining geometric contacts between these components.
Every required contact belongs to the scanned set, so each final
component contains a top and $k\leq h_0$.  Every extra pivot overlap
costs $\varepsilon^3\leq\varepsilon^2$; every other pin has the
$C_T\sqrt{R_1p}\varepsilon^2$ bound just proved.

It remains to integrate the $k$ free translations.  In a final
component all of its source positions have the form $a+d_i$ and
its source velocities are independent of $a$.  Choose one source
label in that component.  Integrate its $g_i(a+d_i)$ using its
$L^1$ norm and bound the other spatial envelopes by their $L^\infty$
norms.  Both norms are at most $C M$ by
\eqref{eq:hist-source-envelope}.  On the original stratum the
restricted component flows equal the full flows; extending these
restricted functions after discarding the stratum indicators
preserves their translation property.  Thus all free translations
cost $C^p$.  The mixed Gaussian integrates to $C^p$, and the speed
powers cost $p^{Cp}$.  This proves the sharper stratum bound and hence
\eqref{eq:hist-pivot-bound}.
\end{proof}

\begin{proposition}[Extensive simultaneous crowding]
\label{hist:crowding}
Fix $A,Q,L,R$ and $\delta>0$.  Among independently enveloped finite
histories with $p\leq P^A$ and $S\leq Rp$, the activity-normalized integral
for which at some time at least $\delta p$ displayed labels have a
companion within $\rho=\varepsilon^{.85}$ is bounded by
\begin{equation}
 \eta_\Gamma P^{C_{A,Q,L,R}p}
       \varepsilon^{.275\delta p-1}.
 \label{eq:hist-crowding-explicit}
\end{equation}
For $\delta p\geq8$ this is at most
$\eta_\Gamma\varepsilon^{.1\delta p}$ for sufficiently small
$\varepsilon$, uniformly in the stated range of $p$.
\end{proposition}

\begin{proof}
The proximity graph of $m$ crowded labels has a forest of at least
$m/2$ edges: each nontrivial connected component with $s$ vertices
contributes $s-1\geq s/2$.  Summing possible forests costs at most
$p^{Cp}$.  Take a grid of spacing at most $\varepsilon$ in $[0,T]$.
On $S\leq Rp$, the displacement of an extended path to its closest
grid point is at most $C_R\sqrt p\,\varepsilon\leq\rho/2$ for
small $\varepsilon$.  Thus an active-path crowding event gives the
same forest with radius $2\rho$ at a grid point, even across a
lifetime boundary.  The grid costs at most $C_T\varepsilon^{-1}$.

Use the sharper form of Lemma~\ref{hist:pivot}.  On a stratum with
$k\leq h_0$, multiplying by
$\mu^{p-h_0}=\varepsilon^{-2(p-h_0)}$ leaves the exponent
\[
 3(.85)r+2(p-r-k)-2(p-h_0)
       =.55r+2(h_0-k)\geq .55r.
\]
Since $r\geq\delta p/2$, this gives
\eqref{eq:hist-crowding-explicit} for a fixed prescription up to a
$p^{Cp}$ factor.  There are only $O_L(p)$ nontrivial cell and
retained-slot choices: empty stretches introduce no new prescription.
Even the crude count $(LJ)^{Cp}p^{C_Lp}$ of all their locations and
finite partitions is $P^{C_{A,Q,L}p}$.  This proves the asserted
summed estimate.  When $\delta p\geq8$, the grid loss leaves an
exponent at least $.15\delta p$; the remaining logarithmic powers
are absorbed by $.05\delta p$ for small $\varepsilon$.
\end{proof}

\begin{corollary}[Separated original parents]
\label{hist:separated}
Given $L,A,\delta>0$ and $M_0<\infty$, one can discard finite terms
of total integral at most $\eta_\Gamma e^{-M_0p}$ so that, except at
fewer than $\delta p$ ordinary pins, the original parent is at
distance greater than $\varepsilon^{.85}$ from every other already
displayed path immediately before the insertion.  This assertion is
used for large histories; fixed bounded orders are treated by their
direct geometric estimates.
\end{corollary}

\begin{proof}
First remove $S>Rp$ with Proposition~\ref{hist:basic}.  Apply
Proposition~\ref{hist:crowding} with a sufficiently small fixed
crowding fraction $\theta$, to be chosen below.  On its complement,
at every ordinary pin at most $\theta p$ of the known old paths
have another known path within $\varepsilon^{.85}$.  Their known
states at the original first-contact time equal the corresponding
states of the prescribed term before the insertion.  Exclude the
new label itself in this count.

In the reference measure retain the adapted gate that the known
old-path crowding count is at most $\theta p$.  It contains every
term under consideration.  Uniform parent selection now gives
conditional probability at most $\theta$ of choosing a crowded old
parent.  Equation~\eqref{eq:hist-reference-discard} bounds the
integral with at least $\delta p$ such choices by
$C_{L,R}^p2^{p/2}\theta^{\delta p/2}$.  Choose $\theta$ small enough
to make this $e^{-(M_0+1)p}$.  The extensive crowding estimate with
this fixed $\theta$ is smaller than the same quantity once the
history size and $P$ are sufficiently large.  This proves the
claim.  Neither step presupposes a good event for the exact cloud.
\end{proof}

All constants used for the last two propositions are chosen after
$L$ and the requested discard exponent.  In contrast, the source
exponent, the parent inequalities, and the constants in
\eqref{eq:hist-finite-bound} were fixed before those choices.  This
separation is what permits the few-crossing estimate in
Section~\ref{sec:closure} to retain its base $Cb\sqrt L$.
